\section{Related Works}
\paragraph{Video Diffusion Models} The field of video generation was revolutionized by the introduction of diffusion models, with pioneering works like AnimateDiff~\cite{guo2023animatediff} and Video Diffusion Models~\cite{ho2022video} marking the beginning of this new era. This initial wave of research was further advanced by models such as Stable Video Diffusion~\cite{blattmann2023stable}, which first demonstrated that curating large-scale, high-quality datasets is critical for enhancing model performance. A common characteristic of these early models was their reliance on UNet architectures. To mitigate the challenges posed by the limited availability and quality of video data compared to images, these models were typically fine-tuned from pre-trained UNet image foundations.
A significant paradigm shift occurred with the introduction of Sora~\cite{openaisora2024}, which heralded a new age defined by the Diffusion Transformer (DiT) architecture~\cite{peebles2023scalable} and training on massive, high-quality video corpora. This breakthrough spurred a proliferation of subsequent research. CogVideox~\cite{yang2024cogvideox} was the first to release an open-source DiT-based model, providing a significant catalyst for community-driven innovation. This was followed by other notable open-source models such as HunyuanVideo~\cite{kong2024hunyuanvideo}, WAN2.1~\cite{wan2025wan}, and Step-Video~\cite{ma2025step}, as well as high-performing closed-source systems including Kling~\cite{kuaishou2024kling}, SeeDance 1.0~\cite{gao2025seedance}, Movie Gen~\cite{polyak2024movie}, and Veo2~\cite{veo22025}.
Architecturally, these contemporary models converge on a common blueprint. They employ a 3D Variational Autoencoder (VAE) to achieve spatio-temporal compression of video into a latent space. The core generative process is then handled by a DiT, which learns to denoise these noisy latents. Across these state-of-the-art models, the quality and scale of the training data have been identified as paramount factors, making data curation and processing a central component of their development.

\paragraph{Audio Diffusion Models} The application of diffusion models to audio generation has followed a parallel trajectory to their video counterparts, fundamentally transforming the landscape of text-to-audio and text-to-music synthesis. Early explorations demonstrated the potential of diffusion for generating high-fidelity audio, but a pivotal advancement was the adoption of latent diffusion architectures~\cite{rombach2022high}, which significantly improved both efficiency and quality. A common technical pipeline for these approaches involves first transforming the raw audio waveform into a mel spectrogram. A Variational Autoencoder (VAE) is then trained on this spectrogram representation to learn a compressed latent space, within which the core diffusion process operates. Within this framework, models such as Stable Audio Open~\cite{evans2025stable} and Riffusion~\cite{riffusion} excel at generating high-fidelity, long-form audio. Furthermore, models like the AudioLDM series~\cite{audioldm2-2024taslp,liu2023audioldm} and DiffRhythm~\cite{ning2025diffrhythm} advance these capabilities to vocal music synthesis, offering fine-grained control over rhythm and other expressive attributes.

\paragraph{Joint Audio and Video Generation} The exploration of joint audio-video generation within diffusion frameworks began with pioneering efforts like MM-Diffusion~\cite{ruan2023mm}. This model was the first to tackle this bimodal task, employing a UNet architecture with two distinct subnetworks, each dedicated to processing the audio and video modalities, respectively. Following this initial work, a series of subsequent models emerged~\cite{ishii2024simple,hayakawa2024mmdisco,ergasti2025r}. However, these early approaches were typically constrained by small-scale training datasets~\cite{lee2022sound,li2021ai}, often less than $10$ hours in size, which inherently limited their diversity and generalization capabilities. A notable step forward was made by models such as Syncflow~\cite{liu2024syncflow} and Uniform~\cite{zhao2025uniform}, which scaled up the training data to approximately $500$ hours by leveraging the VGGSound~\cite{chen2020vggsound} and AudioSet~\cite{gemmeke2017audio} dataset, thereby enhancing their generalization. Despite this progress, persistent challenges remained, including suboptimal video quality and a low degree of disentanglement between the audio and visual streams. The advent of Google's Veo3~\cite{veo32025} marked a significant milestone, representing the first large-scale initiative in synchronous audio-video generation. Veo3 demonstrated the capacity for generating high-fidelity audio and video that is not only diverse but also semantically and temporally coordinated, strictly adhering to user-provided text prompts.